\documentclass[10pt,a4paper]{amsart}
\usepackage[margin=19mm]{geometry}
\usepackage{amsmath,amssymb,mathtools}
\usepackage[scr=rsfs,cal=euler]{mathalfa}
\usepackage{xfrac}
\usepackage[unicode,hidelinks,bookmarks=false]{hyperref}
\newcommand{\ie}{i.e.\ }
\newcommand{\cf}{cf.\ }
\newcommand{\resp}{respectively\ }
\newcommand{\Subsection}{\subsection}
\providecommand{\nobreakdash}{\nobreak\hbox{-}}
\setlength{\emergencystretch}{2em}
\multlinegap0pt
\usepackage{amssymb}
\usepackage[shortlabels]{enumitem}
\usepackage{xcolor}
\usepackage[normalem]{ulem}
\newtheorem{lemma}{Lemma}
\title{Arithmetic Properties of Overcolored Odd Partitions}
\author{M. P. Thejitha, S. N. Fathima}
\date{Source: arXiv:2605.21321v1 (CC BY 4.0)}
\begin{document}
\maketitle

\noindent\emph{Source and licence.} Arithmetic Properties of Overcolored Odd Partitions. Version arXiv:2605.21321v1 (CC BY 4.0), distributed by arXiv under the Creative Commons Attribution 4.0 International licence, http://creativecommons.org/licenses/by/4.0/, as recorded in the arXiv licence metadata for this exact version and shown on the abstract page https://arxiv.org/abs/2605.21321v1. Full licence notice: \texttt{licenses/arxiv-2605.21321-CC-BY-4.0.txt}.\par\smallskip

\noindent\textit{Editorial scope.} Counted unit: the article's lemma lgf2 (source0.tex lines 534--545). The article's complete original proof of it is delivered with it (source0.tex lines 546--643, one \texttt{proof} environment of 98 lines). No mathematics is written, repaired or re-derived: the statement and the proof are the article's own lines, in the article's own notation, and no retained line is substituted or re-flowed.  Retained uncounted as the source-local context the proof needs: lines 148--150; lines 436--441; lines 446--450.  Nothing was omitted from any retained span, so the package carries no omission record.  No retained line contains a citation, so the wrapper carries no bibliography and no bibliography entry.  No retained line contains a figure environment, an image-inclusion command or drawing code, so no image is delivered and the package declares no asset dependency.  Editorial formatting only, in addition: an AMS \texttt{amsart} preamble that restates the article's own declarations and macros used by the retained lines, namely the environments \texttt{lemma} on separate counters, together with the standard packages those lines need.  The retained environments print in this document as Lemma 1 in order of appearance, whereas the article prints the same items with the numbering of its own full document.  The complete native source of the article as distributed in the arXiv e-print for this exact version is bundled unchanged as \texttt{sources/201056/source0.tex}.
	\begin{align}\label{gf}
		\sum_{n=0}^{\infty}\bar{a}_s(n)q^n=\dfrac{f_2^{3s-2}}{f_1^{2s}f_4^{s-1}}.
	\end{align}

\begin{lemma}\label{bt}
	For any prime $p$ and positive integers $k$ and $m$, we have
	\begin{align*}
		f_m^{p^k} \equiv f_{mp}^{p^{k-1}} \pmod{p^k}.
	\end{align*}
\end{lemma}

	\begin{align}
		\dfrac{1}{f_1^2}&=\dfrac{f_8^5}{f_2^5f_{16}^2}+2q\dfrac{f_4^2f_{16}^2}{f_2^5f_8} \label{edf2}\\
		\dfrac{1}{f_1^4}&=\dfrac{f_4^{14}}{f_2^{14}f_8^4}+4q\dfrac{f_4^2f_8^4}{f_2^{10}} \label{edf4}\\
		\dfrac{1}{f_1^8}&=\dfrac{f_4^{28}}{f_2^{28}f_8^8}+8q\dfrac{f_4^{16}}{f_2^{24}}+16q^2 \dfrac{f_4^4f_8^8}{f_2^{20}} \label{edf8}.	
	\end{align}

	\begin{lemma}\label{lgf2}
		For all $n\ge0$ and $\alpha\ge 1$, we have
	\begin{align}
		\sum_{n=0}^{\infty}\bar{a}_{2\alpha+1}(2n+1)q^n&\equiv 2f_1^{12}\pmod4\label{e2no}\\
		\sum_{n=0}^{\infty}\bar{a}_{2\alpha+1}(4n+1)q^n&\equiv 2f_1^{6}\pmod4\label{e4no}\\
		\sum_{n=0}^{\infty}\bar{a}_{2\alpha+1}(4n+2)q^n&\equiv 4f_1^{12}\pmod8\label{e4n2o}\\
		\sum_{n=0}^{\infty}\bar{a}_{2\alpha+1}(4n+3)q^n&\equiv 8f_1^{18}\pmod{16}\label{e4n3o}\\
		\sum_{n=0}^{\infty}\bar{a}_{2\alpha+1}(8n+4)q^n&\equiv 2f_1^{12}\pmod4\label{e8n4o}\\
		\sum_{n=0}^{\infty}\bar{a}_{2\alpha+1}(8n+5)q^n&\equiv 8f_1^{15}\pmod{16}\label{e8n5o}\\
	\sum_{n=0}^{\infty}\bar{a}_{2\alpha+1}(8n+7)q^n&\equiv 64f_1^{21}\pmod{128}\label{e8n7o}.
	\end{align}
\end{lemma}

	\begin{proof}
	Thanks to \eqref{gf}, we have
	\begin{align*}
	\sum_{n=0}^{\infty}\bar{a}_{2\alpha+1}(n)q^n=\dfrac{f_2^{6\alpha+1}}{f_1^{4\alpha+2}f_4^{2\alpha}}.
	\end{align*}
	Employing \eqref{edf2} and \eqref{edf4} in the above equation, we obtain
	\begin{align}
		\sum_{n=0}^{\infty}\bar{a}_{2\alpha+1}(n)q^n&=\dfrac{f_2^{6\alpha+1}}{f_4^{2\alpha}}\left(\dfrac{f_8^5}{f_2^5f_{16}^2}+2q\dfrac{f_4^2f_{16}^2}{f_2^5f_8}\right)\left(\dfrac{f_4^{14}}{f_2^{14}f_8^4}+4q\dfrac{f_4^2f_8^4}{f_2^{10}}\right)^\alpha\nonumber\\
		&=\dfrac{f_2^{6\alpha+1}}{f_4^{2\alpha}}\left(\dfrac{f_4^{14}}{f_2^{14}f_8^4}\right)^\alpha\left(\dfrac{f_8^5}{f_2^5f_{16}^2}+2q\dfrac{f_4^2f_{16}^2}{f_2^5f_8}\right)\left(\sum_{i=0}^{\infty}\binom{\alpha}{i}4^iq^i\dfrac{f_2^{4i}f_8^{8i}}{f_4^{12i}}\right)\nonumber
	\end{align}
	\begin{align}
		&=\dfrac{f_4^{12\alpha}}{f_2^{8\alpha-1}f_8^{4\alpha}}\left(\dfrac{f_8^5}{f_2^5f_{16}^2}\sum_{i=0}^{\alpha}\binom{\alpha}{i}4^iq^i\dfrac{f_2^{4i}f_8^{8i}}{f_4^{12i}}+2\dfrac{f_4^2f_{16}^2}{f_2^5f_8}\sum_{i=0}^{\alpha}\binom{\alpha}{i}4^iq^{i+1}\dfrac{f_2^{4i}f_8^{8i}}{f_4^{12i}}\right)\label{el8}
	\end{align}
		Extracting the terms of the form $q^{2n+1}$ from both sides of \eqref{el8}, and then replacing $q^{2}$ with $q$ gives
	\begin{align}
	\sum_{n=0}^{\infty}\bar{a}_{2\alpha+1}(2n+1)q^n&=\dfrac{f_2^{12\alpha}}{f_1^{8\alpha-1}f_4^{4\alpha}}\left(\dfrac{f_4^5}{f_1^5f_{8}^2}\sum_{i=0}^{\alpha}\binom{\alpha}{2i+1}4^{2i+1}q^i\dfrac{f_1^{8i+4}f_4^{16i+8}}{f_2^{24i+12}}\right.\nonumber\\
	&\left.+2\dfrac{f_2^2f_{8}^2}{f_1^5f_4}\sum_{i=0}^{\alpha}\binom{\alpha}{2i}4^{2i}q^{i}\dfrac{f_1^{8i}f_4^{16i}}{f_2^{24i}}\right)\nonumber\\
	&\equiv 2\dfrac{f_2^{12\alpha+2}f_8^2}{f_1^{8\alpha+4}f_4^{4\alpha+1}}\pmod 4.\label{el4}
	\end{align}
	Employing Lemma \ref{bt} in the above congruence yields \eqref{e2no}.\\
\indent	Using \eqref{edf4} and \eqref{edf8} in \eqref{el4}, we have
\begin{align}
\sum_{n=0}^{\infty}\bar{a}_{2\alpha+1}(2n+1)q^n&\equiv 2\dfrac{f_2^{12\alpha+2}f_8^2}{f_4^{4\alpha+1}}\left(\dfrac{f_4^{14}}{f_2^{14}f_8^4}+4q\dfrac{f_4^2f_8^4}{f_2^{10}}\right)\nonumber\\
&\cdot\left(\dfrac{f_4^{28}}{f_2^{28}f_8^8}+8q\dfrac{f_4^{16}}{f_2^{24}}+16q^2 \dfrac{f_4^4f_8^8}{f_2^{20}}\right)^\alpha\pmod 4\label{el6}\\
&\equiv 2 \dfrac{f_4^{24\alpha+13}}{f_2^{16\alpha+12}f_8^{8\alpha+2}}
\pmod4.\nonumber
\end{align}
	Extracting the terms of the form $q^{2n}$ from both sides of the above equation, and then replacing $q^{2}$ with $q$ yields
\begin{align}\label{el5}
\sum_{n=0}^{\infty}\bar{a}_{2\alpha+1}(4n+1)q^n\equiv 2\dfrac{f_2^{24\alpha+13}}{f_1^{16\alpha+12}f_4^{8\alpha+2}}\pmod 4.
\end{align}
	Employing Lemma \ref{bt} in the above congruence gives \eqref{e4no}.\\
	\indent From \eqref{el4}, we have
	\begin{align*}
	\sum_{n=0}^{\infty}\bar{a}_{2\alpha+1}(2n+1)q^n
	&\equiv 4\alpha\dfrac{f_2^{12\alpha-12}}{f_1^{8\alpha}f_4^{4\alpha-13}f_8^2}+2\dfrac{f_2^{12\alpha+2}f_8^2}{f_1^{8\alpha+4}f_4^{4\alpha+1}}\pmod {16}.
	\end{align*}
	Employing \eqref{edf4} and \eqref{edf8} in the above equation, we get
	\begin{align*}
	\sum_{n=0}^{\infty}\bar{a}_{2\alpha+1}(2n+1)q^n&\equiv 4\alpha\dfrac{f_2^{12\alpha-12}}{f_4^{4\alpha-13}f_8^2}\left(\dfrac{f_4^{28}}{f_2^{28}f_8^8}\right)^\alpha\left(\sum_{i=0}^{\alpha}\binom{\alpha}{i}8^iq^i\dfrac{f_2^{4i}f_8^{8i}}{f_4^{12i}}\right)+2\dfrac{f_2^{12\alpha+2}f_8^2}{f_4^{4\alpha+1}}\\
	&\left(\dfrac{f_4^{28}}{f_2^{28}f_8^{8}}\right)^\alpha\left(\dfrac{f_4^{14}}{f_2^{14}f_8^4}+4q\dfrac{f_4^2f_8^4}{f_2^{10}}\right)\left(\sum_{j=0}^{\alpha}\binom{\alpha}{j}8^jq^j\dfrac{f_2^{4j}f_8^{8j}}{f_4^{12j}}\right)\\
	&\equiv 4\alpha\dfrac{f_4^{24\alpha+13}}{f_2^{16\alpha+12}f_8^{8\alpha+2}}\sum_{i=0}^{\alpha}\binom{\alpha}{i}8^iq^i\dfrac{f_2^{4i}f_8^{8i}}{f_4^{12i}}+2\dfrac{f_4^{24\alpha-1}}{f_2^{16\alpha-2}f_8^{8\alpha-2}}\left(\dfrac{f_4^{14}}{f_2^{14}f_8^4}\right.\\
	&\left.\sum_{j=0}^{\alpha}\binom{\alpha}{j}8^jq^j\dfrac{f_2^{4j}f_8^{8j}}{f_4^{12j}}+4\dfrac{f_4^2f_8^4}{f_2^{10}}\sum_{j=0}^{\alpha}\binom{\alpha}{j}8^jq^{j+1}\dfrac{f_2^{4j}f_8^{8j}}{f_4^{12j}}\right)\pmod{16}.
	\end{align*}
		Extracting the terms of the form $q^{2n+1}$ from both sides of the above equation, and then replacing $q^{2}$ with $q$ yields
	\begin{align}\label{el7}
	\sum_{n=0}^{\infty}\bar{a}_{2\alpha+1}(4n+3)q^n&\equiv8\dfrac{f_2^{24\alpha+1}}{f_1^{16\alpha+8}f_4^{8\alpha-6}} \pmod{16}.
	\end{align}
	Employing Lemma \ref{bt} in the above equation gives \eqref{e4n3o}.\\
\indent From \eqref{el8}, we have
\begin{align*}
\sum_{n=0}^{\infty}\bar{a}_{2\alpha+1}(n)q^n\equiv\dfrac{f_4^{12\alpha}}{f_2^{8\alpha+4}f_8^{4\alpha-5}f_{16}^2}+2q\dfrac{f_4^{12\alpha+2}f_{16}^2}{f_2^{8\alpha+4}f_8^{4\alpha+1}}\pmod 4.
\end{align*}
Extracting the terms of the form $q^{2n}$ from both sides of the above equation, and then replacing $q^{2}$ with $q$ gives
\begin{align}\label{el9}
\sum_{n=0}^{\infty}\bar{a}_{2\alpha+1}(2n)q^n\equiv\dfrac{f_2^{12\alpha}}{f_1^{8\alpha+4}f_4^{4\alpha-5}f_{8}^2}\pmod 4.
\end{align}
Employing \eqref{edf4} and \eqref{edf8} in \eqref{el9}, we obtain
\begin{align}\label{el11}
\sum_{n=0}^{\infty}\bar{a}_{2\alpha+1}(2n)q^n\equiv\dfrac{f_4^{24\alpha+19}}{f_2^{16\alpha+14}f_8^{8\alpha+6}}+4q\dfrac{f_4^{24\alpha+7}}{f_2^{16\alpha+10}f_8^{8\alpha-2}}\pmod 8.
\end{align}
Extracting the terms of the form $q^{2n+1}$ from both sides of the above equation, and then replacing $q^{2}$ with $q$ yields
	\begin{align}\label{el10}
	\sum_{n=0}^{\infty}\bar{a}_{2\alpha+1}(4n+2)q^n\equiv4\dfrac{f_2^{24\alpha+7}}{f_1^{16\alpha+10}f_4^{8\alpha-2}}\pmod 8.
	\end{align}
	Employing Lemma \ref{bt} in \eqref{el10} gives \eqref{e4n2o}.\\
	\indent	Extracting the terms of the form $q^{2n}$ from both sides of \eqref{el11}, and then replacing $q^{2}$ with $q$ yields 
		\begin{align*}
		\sum_{n=0}^{\infty}\bar{a}_{2\alpha+1}(4n)q^n\equiv\dfrac{f_2^{24\alpha+19}}{f_1^{16\alpha+14}f_4^{8\alpha+6}}\pmod 4.
		\end{align*}
			Employing \eqref{edf2} in the above equation, we obtain
		\begin{align*}
			\sum_{n=0}^{\infty}\bar{a}_{2\alpha+1}(4n)q^n\equiv\dfrac{f_2^{24\alpha+19}}{f_4^{8\alpha+6}}\left(\dfrac{f_8^5}{f_2^5f_{16}^2}\right)^{8\alpha+7}\sum_{i=0}^{8\alpha+7}\binom{8\alpha+7}{i}2^iq^i\dfrac{f_4^{2i}f_{16}^{4i}}{f_8^{6i}}\pmod 4.
		\end{align*}
		Extracting the terms of the form $q^{2n+1}$ from both sides of the above equation, and then replacing $q^{2}$ with $q$ yields
			\begin{align*}
			\sum_{n=0}^{\infty}\bar{a}_{2\alpha+1}(8n+4)q^n\equiv 2\dfrac{f_4^{40\alpha+29}}{f_1^{16\alpha+16}f_2^{8\alpha+4}f_8^{16\alpha+10}}\pmod 4.
			\end{align*}
			Employing Lemma \ref{bt} in the above equation gives \eqref{e8n4o}.\\
		\indent	Employing \eqref{edf2} in \eqref{el5}, we have
			\begin{align*}
			\sum_{n=0}^{\infty}\bar{a}_{2\alpha+1}(4n+1)q^n&\equiv 2\dfrac{f_2^{24\alpha+13}}{f_4^{8\alpha+2}}\left(\dfrac{f_8^5}{f_2^5f_{16}^2}\right)^{8\alpha+6}\sum_{i=0}^{8\alpha+6}\binom{8\alpha+6}{i}2^iq^i\dfrac{f_4^{2i}f_{16}^{4i}}{f_8^{6i}}\pmod 4.
			\end{align*}
			Extracting the terms of the form $q^{2n+1}$ from both sides of the above equation, and then replacing $q^{2}$ with $q$ gives
			\begin{align*}
			\sum_{n=0}^{\infty}\bar{a}_{2\alpha+1}(8n+5)q^n\equiv 8 \dfrac{f_4^{40\alpha+24}}{f_1^{16\alpha+17}f_2^{8\alpha}f_8^{16\alpha+8}}\pmod{16}.
			\end{align*}
			Employing Lemma \ref{bt} in the above equation results in \eqref{e8n5o}.\\
			\indent Employing \eqref{edf2} in \eqref{el7}, we have
			\begin{align*}
			\sum_{n=0}^{\infty}\bar{a}_{2\alpha+1}(4n+3)q^n\equiv 8\dfrac{f_2^{24\alpha+1}}{f_4^{8\alpha-6}}\left(\dfrac{f_8^5}{f_2^5f_{16}^2}\right)^{8\alpha+4}\sum_{i=0}^{8\alpha+4}\binom{8\alpha+4}{i}2^iq^i\dfrac{f_4^{2i}f_{16}^{4i}}{f_8^{6i}}\pmod{16}.
			\end{align*}
			Extracting the terms of the form $q^{2n+1}$ from both sides of the above equation, and then replacing $q^{2}$ with $q$ yields
			\begin{align*}
			\sum_{n=0}^{\infty}\bar{a}_{2\alpha+1}(8n+7)q^n\equiv 64\dfrac{f_4^{40\alpha+14}}{f_1^{16\alpha+19}f_2^{8\alpha-8}f_8^{16\alpha+4}}\pmod{128}.
			\end{align*}
			Employing Lemma \ref{bt} in the above equation, we arrive at \eqref{e8n7o}.
	\end{proof}

\end{document}
